\documentclass[10pt]{article}
\usepackage{graphicx}
\usepackage{lipsum}  % For sample text
\usepackage{url}
\usepackage{booktabs} % For formal tables
\usepackage{amsmath} % For better math formatting
\usepackage{geometry} % For better margin settings
\usepackage{titlesec} % For customizing sections and subsections
\usepackage{siunitx}
\usepackage[colorlinks=true, allcolors=black, linktoc=all]{hyperref}
\usepackage{tikz}
\usepackage{fancyhdr}
\usetikzlibrary{shapes.geometric, arrows}

\usepackage{algorithm}
\usepackage{algpseudocode}
\usepackage{amsmath}
\usepackage{listings}
\usepackage{float}
\usepackage{hyperref}


\lstset{
    basicstyle=\ttfamily,
    captionpos=b,
    frame=single,
    numbers=left,
    breaklines=true,
}

\pagestyle{fancy}
\fancyhf{} 
\lhead{\textit{ECE 652}} 
\rhead{\thepage}  

% Customize section and subsection formatting
\titleformat{\section}
{\normalfont\Large}{\thesection}{1em}{}
\titleformat{\subsection}
{\normalfont\large}{\thesubsection}{1em}{}
\titleformat{\subsubsection}
{\normalfont}{\thesubsubsection}{1em}{}

% Document title and author
\title{FPGA Implementation of Coherent Multicore Processor}
\author{
    Roshan Sundaram \\
    \texttt{roshan.sundaram@duke.edu}
}
\date{\today}

\begin{document}

\maketitle

\begin{center}    


% Abstract will be in center \\ Remove center to make it standard

\end{center}


\section{Introduction}
Given that ECE 652 is a course focused on multicore architectures and the structures that they are built upon, this project will implement the topics discussed in class in hardware. In ECE 350 (digital design), we built a 5-stage pipelined CPU in Verilog that implemented a modified version of the MIPS architecture. Taking inspiration from that endeavor, this project will focus on integrating multiple CPU cores in Verilog and tying them together with a simple but functional cache coherence policy. The multicore processor will then be synthesized and implemented on an FPGA. 

\section{Background}
Multicore processors make up the vast majority of modern chip designs. However, connecting multiple cores together and exploiting thread-level parallelism relies on the fundamental assumption that the memory system will behave in a predictable, predefined manner. This requirement is known as memory consistency, and is concerned with the global ordering of load and store instructions. Consistency, being an architectural specification, is visible to the programmer. \textit{Micro}architects may implement various microarchitectural features to guarantee consistency, however. To that end, cache coherence protocols are systems that support consistency models but are invisible to the programmer. Various coherence protocols -- or no coherence at all -- may be used to uphold a consistency model. Within cache coherence, two popular solutions exist: snooping protocols, and directory protocols. 

\textbf{Snooping protocols} require cores to broadcast coherence messages regarding the state of a certain cache block to a shared, totally-ordered bus. Each cache monitors the bus and, upon receiving a relevant coherence message from another core, performs the required action on its own cache. Such protocols are conceptually simple since they abstract away the complexities and challenges associated with a ``shared bus'' but are difficult to scale. 

\textbf{Directory protocols} remove some of the overhead of snooping protocols by adding a data structure known as the directory to the last-level cache. For each cache block, the directory tracks which cores have the block, the state of each block in each core, and the owner of the block. When transactions are issued, instead of the message being put onto a shared bus, it is sent directly to the directory, at which point the transaction becomes ordered and is serviced. 

The states tracked by a coherence protocol can be as few as two and up to a very large number (if counting transient states). The simplest protocol (applied generally to a snooping or directory protocol) is a valid/invalid protocol, which tracks whether a block is valid or invalid using a validity bit. While such a system is easy to implement, it leads to excessive overhead from coherence transactions. More effective protocols are those which are a subset of the  ``MOESI'' states. In these protocols, each block is said to be in one of five (or fewer) stable states representing the ownership and sharing information of that cache block. When cores read or write data, the protocol dictates state changes. 

Connecting multiple cores together in hardware is the job of the interconnection network, which can take on a variety of topologies. Within a physical network, a designer can implement virtual networks, virtual channels, and buffers at the source and destination ends. A robust network is crucial to preventing deadlocks (protocol and routing) and ensuring smooth operation of the multicore processor. For a simple multiprocessor, a simple network may suffice: rings and meshes are conceptually simple and can achieve decent performance. 


\section{Implementation and Execution}
\subsection{Hardware Description Language}

This project will be designed and implemented in a hardware description language (HDL) and synthesized in AMD Vivado following steps similar to those practiced in ECE 350. The basis of the processor will be the WARP-V configurable RISC core developed by Steve Hoover of Redwood EDA \cite{warpv}. The WARP-V configurator allows for the generation of a pipelined processor (RISC-V or MIPS) in transaction-level Verilog (TL-Verilog) \cite{tlx}. TL-Verilog is an HDL built upon SystemVerilog that allows for easy design of pipelined CPUs without explicitly declaring latches and wires. TL-Verilog is then translated to SystemVerilog using Redwood's SandPiper tool, which can be accessed online through their MakerChip development tool or run locally. \footnote{I have requested an educational copy of SandPiper so I can run the software locally; however, if my request is not approved, the online version should suffice.} SystemVerilog can then be synthesized in Vivado and deployed to an FPGA. \footnote{The exact model of FPGA is yet to be determined and will depend on the number of logic cells needed.}

\subsection{Cores}
WARP-V provides a detailed configurator that can generate everything from a simple, single-cycle processor to a deeply-pipelined, superscalar core. For this project, two simple cores will be used, likely with three to five pipeline stages each for understandability. The cores themselves are not generated with any IO capabilities or cache. Thus, it will be a major effort of this project to implement the private caches and perhaps a shared L2 cache. Multi-cycle functional units like multipliers and dividers may also be omitted due to their high chip area. 

\subsection{Interconnection Network} 
Given that this project will consist of just two cores to start off, the choice of interconnection network becomes trivial as scaling will not be an issue. Each core must be linked to the other and to the last level cache (LLC). This goal could be achieved via three bidirectional links (C1 to C2, C1 to LLC, and C2 to LLC). This framework could be modified easily to implement either a snooping or directory protocol. 

\subsection{Coherence Protocol}
A snooping protocol will be used to implement cache coherence in this multicore processor, as it is conceptually simple and only requires the addition of state bits to each cache block and the construction of a bus arbiter to handle interactions between the bus and the cores/caches. This module will be constructed in TL-Verilog or SystemVerilog (depending on the capabilities of TL-Verilog) and will contain the bulk of the logic of the project. 

An MSI coherence protocol is a robust protocol that will be used in this project. Although a valid/invalid protocol would have been simpler, MSI offers a better balance of complexity and performance, and can be implemented by following the state transition diagram discussed in class. The analysis of transient states will be a challenge but existing research and literature on MSI protocols will be studied and, if needed, a model checker like Murphi can be used to verify the protocol. 



\section{Experimental Setup}
Testing at every step is crucial to a successful project. First, each Verilog component will be tested using testbenches to verify correct behavior. The skills learned in ECE 350 -- notably the use of GTKWave -- will assist in the debugging effort. The testbenches will emulate surrounding hardware -- such as a dummy core to test the arbiter, or a dummy bus to test the L1 cache -- to rigorously verify behavior under a variety of scenarios Once each component is verified, the full system will be built from the various individual components and tested using microbenchmarks. Microbenchmarks will be written in C, compiled down to either RISC-V or MIPS assembly, and run on the multicore processor. \footnote{I still have not decided between RISC-V and MIPS. The WARP-V configurator supports both, but I am more familiar with MIPS.} Once the final system is synthesized and implemented on the FPGA, the integrated logic analyzer (ILA) and debug cores in Vivado will be used to observe the operation of the coherence protocol and compare against expected values. This stage of testing will be the most time-consuming, as the ILA is notoriously slow to synthesize and use. 

\section{Conclusion}
The goal of this project is less to advance a cutting-edge area of research and more to implement some very interesting topics in hardware and get practice with building real CPUs. A reasonable timeline for this project would be to generate the WARP-V cores, design the bus \& caches in Verilog, and synthesize a basic version of the multicore in Vivado over spring break. Then, during the remainder of March, the custom bus arbiter will be designed in Verilog to implement the MSI protocol. Finally, the remaining weeks in April will see extensive testing of the full system and the writing of the final paper.


\clearpage 


\begin{thebibliography}{99}

\bibitem{warpv}
Hoover, Steve. "WARP-V." \textit{GitHub Repository}.
\url{https://github.com/stevehoover/warp-v}

\bibitem{tlx}
"TL-X." \textit{Transaction-Level Design}.
\url{https://www.tl-x.org}

\end{thebibliography}





\end{document}